\section{Complete modular reconstruction for \texorpdfstring{$\RR{4}{155}$}{R6\_4\_155}}\label{sec:155-complete-modular}

Write $d=4+3\sqrt2$, $D=72+48\sqrt2$, and $z=e^{2\pi i/24}$.
Section~\ref{sec:frontier-155} gives
\[
 Z(\mathcal C)\simeq\mathcal S\boxtimes\mathcal B,
 \qquad \operatorname{rank}\mathcal B=26,\qquad
 \dim\mathcal B=144d^2,
\]
where $\mathcal S$ is the semion category. The multiplicative central
charge of $\mathcal B$ is $z^{-3}$. Thus
\[
 s=\frac{S_{\mathcal B}}{12d},\qquad t=zT_{\mathcal B}
\]
obey $s^2=C$ and $(st)^3=C$, where $C$ is charge conjugation.
Dong--Lin--Ng's congruence and Galois theorem gives a representation of
$\mathrm{SL}_2(\mathbb Z/24\mathbb Z)$ with entries in
$\mathbb Q(z)$; see \href{https://arxiv.org/abs/1201.6644}{arXiv:1201.6644},
Theorem II and Section 4. All labels in this section are the zero-based
labels of the frozen rank-26 restriction list; the vacuum is label $1$.

\subsection{Two inexpensive duality obstructions}

\begin{lemma}[Determinant obstruction]
If a modular representation has $s^2=C$ and $(st)^3=C$, then
$\det C=(\det t)^{-6}$.
For the retained spectrum of $\mathcal B$, the number $r$ of self-dual
simples satisfies $r\equiv26\pmod4$.
\end{lemma}
\begin{proof}
Taking determinants gives $\det s=(\det t)^{-3}$.
The twist exponents in $T_{\mathcal B}$ sum to $258$, so
$\det t=z^{258+26}=z^{20}$ and $\det C=1$.
An involution on 26 labels with $r$ fixed labels has determinant
$(-1)^{(26-r)/2}$. The previously allowed value $r=12$ is therefore
impossible; the two remaining values are $10$ and $14$.
\end{proof}

For every $a$ coprime to 24, $a^2=1\pmod{24}$. Hence the signed Galois
matrix is
\[
 G_a=\sigma_a(s)s^{-1}=\rho(aI).
\]
It is central, involutory, and commutes with $s,t$. Its simultaneous
eigenspaces have rational bases. If $G_a$ acts on such a space by a
character $\chi(a)$, each entry of its $s$ block transforms by that
character. Write $\chi_i,\chi_2,\chi_3$ for the characters belonging
to $i,\sqrt2,\sqrt3$.

The exact character-table calculation described below forces a
12-dimensional trivial-character block, a six-dimensional $\chi_2$
block, a six-dimensional $\chi_i$ block, and a final two-dimensional
block. That last block has character $\chi_i$ when the two odd
$2d$ pairs are dual pairs, and $\chi_3$ when both pairs are self-dual.

\begin{lemma}[Rational-square obstruction]
The self-dual odd-$2d$ branch is impossible. Consequently $\mathcal B$
has exactly ten self-dual simples.
\end{lemma}
\begin{proof}
In the self-dual branch use the equal-norm rational basis
$e_{12}-e_{13},e_{14}-e_{15}$. On this $\chi_3$ space,
\[
 t=\operatorname{diag}(i,iw),\qquad
 s=\begin{pmatrix}a&b\\b&c\end{pmatrix},\qquad w=e^{2\pi i/3},
\]
and $a,b,c\in\sqrt3\,\mathbb Q$. The block is irreducible, so
$b\ne0$. The relation $s^2=I$ gives $c=-a$.
The off-diagonal entry of $sts=t^{-1}st^{-1}$ gives
\[
 a=\frac{1}{i(iw)(i-iw)}=\frac1{\sqrt3},\qquad b^2=\frac23.
\]
Thus $(b/\sqrt3)^2=2/9$, impossible for a rational number.
This is a field obstruction, not a numerical approximation.
\end{proof}

\subsection{Exact determination of every $S$ entry}

The replay first constructs and verifies the class algebras and
character tables of $\mathrm{SL}_2(\mathbb Z/8\mathbb Z)$ and
$\mathrm{SL}_2(\mathbb Z/3\mathbb Z)$. Their 210 tensor-product
irreducibles exhaust the irreducibles at level 24. There are 718
dimension-26 decompositions with the prescribed $t$ spectrum. Signed
Galois permutations and the known nonreal-character rows leave eight.
The known $\chi_2$ block has $\operatorname{Tr}(st)=0$ and leaves four.
The preceding rational-square obstruction removes the self-dual
branch.

On the mandatory $\chi_i$ six-space, take the pair-difference basis
\[
(2,5),\ (3,6),\ (4,7),\ (20,21),\ (16,17),\ (18,19),
\]
where $(a,b)$ denotes $e_a-e_b$. The block is
\[
 s=\frac i2
 \begin{pmatrix}
0&0&1&-1&-1&-1\\0&0&-1&1&-1&-1\\
1&-1&0&0&1&-1\\-1&1&0&0&1&-1\\
-1&-1&1&1&0&0\\-1&-1&-1&-1&0&0
 \end{pmatrix}.
\]
The last two rows are known low rows. Their modular relations leave
three rational variables; their square relations force them to
$0,1/2,0$. The additional $\chi_i$ two-space has $s=iI$.

Only the rational 12-dimensional block remains. Its basis is
\[
\begin{gathered}
e_1-e_{23},\ e_0-e_{22},\ e_2+e_5,\ e_3+e_6,\ e_4+e_7,
\ e_{20}+e_{21},\\
e_8,\ e_9,\ e_{10},\ e_{11},\ e_{12}+e_{13},\ e_{14}+e_{15}.
\end{gathered}
\]
Let $P$ have these basis vectors as columns,
$N=P^{\mathsf T}P=\operatorname{diag}(2,2,2,2,2,2,1,1,1,1,2,2)$,
and $H=P^{\mathsf T}sP$. Its known first two rows, together with their
square and modular relations, leave ten rational variables. Two further
linear equations from the full system force one variable to zero.
Expanding the remaining modular relations in the cyclotomic basis then
gives 187 rational polynomials of degree at most two in the nine remaining
variables. Row reduction on their
monomials produces the nine linear consequences
\[
\begin{array}{lll}
h_{7,11}=-1/3,&h_{7,7}=-1/6,&h_{7,8}=1/3,\\
h_{7,9}=1/3,&h_{8,11}=2/3,&h_{8,8}=-1/6,\\
h_{8,9}=-1/6,&h_{9,11}=-2/3,&h_{9,9}=-1/6.
\end{array}
\]
Substitution satisfies all the original equations. Thus the solution
is unique; the final replay uses rational linear elimination, not a
Gr\"obner calculation. The integer matrix $6H$ and all basis conventions
are recorded in \texttt{matrix\_blocks.json}. Its trace is
$\operatorname{Tr}(N^{-1}H)=-2$, selecting the $3+9$ representation
type and excluding the $3+3+3+3$ alternative.

\begin{theorem}[Unique modular candidate]
In the retained normalization and labels, the necessary rank-26
modular datum is unique. Its modular representation has irreducible
dimensions $9,6,4,3,2,1,1$. Every one of its $26^3$ Verlinde
coefficients is a nonnegative integer, with maximum 9.
\end{theorem}
\begin{proof}
The preceding block determinations exhaust the 26-dimensional space.
For each block its contribution to the original normalized matrix is
$PN^{-1}HN^{-1}P^{\mathsf T}$. Their sum is therefore forced.
The exact replay checks every matrix entry of $s^2=C$ and
$sts=t^{-1}st^{-1}$, all known low rows, and every Verlinde coefficient
in $\mathbb Q(z)$. A separate implementation verifies all fusion
associativity identities and exact diagonalization by the reconstructed
$s$, as well as direct Verlinde computations for four selected fusion
matrices. These are checks of the forced candidate, not associators
for a fusion category.
\end{proof}

\subsection{What the completed data pass}

All 456,976 ordered four-object Vafa equations hold exactly modulo 24.
The ordinary second indicators agree with duality: the ten self-dual
simples have indicator $+1$ except label 24, which has indicator $-1$.
It occurs with multiplicity zero in every $U^*\otimes U$ for simple
$U$. The 182 ordinary cyclic spectra for the 26 simples at orders
$2,3,4,6,8,12,24$ have nonnegative integral eigenvalue multiplicities.
All 676 products also agree with the frozen restriction functor to
the original rank-six ring.

Using the full rank-52 matrix for $Z(\mathcal C)$ strengthens the previous
low-target generalized-indicator checks. All 7,176 cyclic spectra for
52 central targets, six original simples and orders 2 through 24
pass, together with 312 degree-one normalization identities.
For $\mathcal B$ itself, all 17,576 generalized second-indicator blocks
and all 17,576 degree-one identities pass.
The trace formula is Ng--Schauenburg,
\href{https://arxiv.org/abs/0806.2493}{arXiv:0806.2493},
Corollary 5.6, with Lemmas 2.5--2.7 and Proposition 2.8.
In particular a target of twist $\theta$ requires the normalization
$E^n=\theta^{-1}$ before cyclic Fourier inversion.
These are finite necessary checks. The exponent 24 alone does not
turn them into a proof of feasibility at every tensor order.

No known category has been identified with these data. Matching
dimensions alone is insufficient, and the new near-group constructions
discussed next do not provide the required condensation.

